% GENERATED FILE. Edit canonical lessons, then run npm run generate:books.
% canonical-source-hash: fnv1a64:4d6580c396ad0586
% canonical-lessons: SA-C188-trsitah, SA-C188-sajjah, SA-C188-vyastah, SA-C188-nirmalah, SA-C188-ujjvalah

\chapter{Thirsty, Ready, Busy, Spotless, Bright}
\label{ch:sa-thirsty-ready-busy-spotless-bright}
\hlchaptermodality{188}

\begin{chapteropening}
\textbf{By the end of this chapter:} \emph{I can say thirsty, ready, busy, spotless and bright.}

Complete the last lesson of chapter 188: I can say thirsty, ready, busy, spotless and bright.
\end{chapteropening}

\section[\texorpdfstring{\sk{तृषितः} (tṛṣitaḥ)}{तृषितः (tṛṣitaḥ)}]{\sk{तृषितः} (tṛṣitaḥ) — thirsty}

\label{lesson:SA-C188-trsitah}



\begin{quote}
Before the new one: say the Sanskrit for truthful, then the Sanskrit for adept.
\end{quote}

\subsection*{You'll want to know: \sk{तृषितः}}
\textbf{\sk{तृषितः}} — \emph{tṛṣitaḥ} — \textquotedblleft{}thirsty\textquotedblright{}.

\begin{grammarlens}[title={What you've built}]
One more word for this chapter.
\end{grammarlens}

\subsection*{Your turn}
\begin{itemize}
\raggedright
  \item \emph{Say it:} \emph{tṛṣitaḥ}
  \item \emph{Say it:} \emph{tṛṣitaḥ}, once more
  \item \emph{Say it:} say \emph{tṛṣitaḥ}
  \item \emph{Recall:} say the Sanskrit for truthful, then the Sanskrit for adept, then say \emph{tṛṣitaḥ} again
\end{itemize}

\begin{tcolorbox}[breakable,colback=teal!4,colframe=teal!35!black,title={Before you move on}]
What does \sk{तृषितः} mean? (\textquotedblleft{}Thirsty\textquotedblright{}.) Say it once more.
\end{tcolorbox}
\section[\texorpdfstring{\sk{सज्जः} (sajjaḥ)}{सज्जः (sajjaḥ)}]{\sk{सज्जः} (sajjaḥ) — ready}

\label{lesson:SA-C188-sajjah}



\begin{quote}
Before the new one: say the Sanskrit for adept, then the Sanskrit for thirsty.
\end{quote}

\subsection*{You'll want to know: \sk{सज्जः}}
\textbf{\sk{सज्जः}} — \emph{sajjaḥ} — \textquotedblleft{}ready\textquotedblright{}.

\begin{grammarlens}[title={What you've built}]
One more word for this chapter.
\end{grammarlens}

\subsection*{Your turn}
\begin{itemize}
\raggedright
  \item \emph{Say it:} \emph{sajjaḥ}
  \item \emph{Say it:} \emph{sajjaḥ}, once more
  \item \emph{Say it:} say \emph{sajjaḥ}
  \item \emph{Recall:} say the Sanskrit for adept, then the Sanskrit for thirsty, then say \emph{sajjaḥ} again
\end{itemize}

\begin{tcolorbox}[breakable,colback=teal!4,colframe=teal!35!black,title={Before you move on}]
What does \sk{सज्जः} mean? (\textquotedblleft{}Ready\textquotedblright{}.) Say it once more.
\end{tcolorbox}
\section[\texorpdfstring{\sk{व्यस्तः} (vyastaḥ)}{व्यस्तः (vyastaḥ)}]{\sk{व्यस्तः} (vyastaḥ) — busy}

\label{lesson:SA-C188-vyastah}



\begin{quote}
Before the new one: say the Sanskrit for thirsty, then the Sanskrit for ready.
\end{quote}

\subsection*{You'll want to know: \sk{व्यस्तः}}
\textbf{\sk{व्यस्तः}} — \emph{vyastaḥ} — \textquotedblleft{}busy\textquotedblright{}.

\begin{grammarlens}[title={What you've built}]
One more word for this chapter.
\end{grammarlens}

\subsection*{Your turn}
\begin{itemize}
\raggedright
  \item \emph{Say it:} \emph{vyastaḥ}
  \item \emph{Say it:} \emph{vyastaḥ}, once more
  \item \emph{Say it:} say \emph{vyastaḥ}
  \item \emph{Recall:} say the Sanskrit for thirsty, then the Sanskrit for ready, then say \emph{vyastaḥ} again
\end{itemize}

\begin{tcolorbox}[breakable,colback=teal!4,colframe=teal!35!black,title={Before you move on}]
What does \sk{व्यस्तः} mean? (\textquotedblleft{}Busy\textquotedblright{}.) Say it once more.
\end{tcolorbox}
\section[\texorpdfstring{\sk{निर्मलः} (nirmalaḥ)}{निर्मलः (nirmalaḥ)}]{\sk{निर्मलः} (nirmalaḥ) — spotless, pure}

\label{lesson:SA-C188-nirmalah}



\begin{quote}
Before the new one: say the Sanskrit for ready, then the Sanskrit for busy.
\end{quote}

\subsection*{You'll want to know: \sk{निर्मलः}}
\textbf{\sk{निर्मलः}} — \emph{nirmalaḥ} — \textquotedblleft{}spotless, pure\textquotedblright{}.

\begin{grammarlens}[title={What you've built}]
One more word for this chapter.
\end{grammarlens}

\subsection*{Your turn}
\begin{itemize}
\raggedright
  \item \emph{Say it:} \emph{nirmalaḥ}
  \item \emph{Say it:} \emph{nirmalaḥ}, once more
  \item \emph{Say it:} say \emph{nirmalaḥ}
  \item \emph{Recall:} say the Sanskrit for ready, then the Sanskrit for busy, then say \emph{nirmalaḥ} again
\end{itemize}

\begin{tcolorbox}[breakable,colback=teal!4,colframe=teal!35!black,title={Before you move on}]
What does \sk{निर्मलः} mean? (\textquotedblleft{}Spotless, pure\textquotedblright{}.) Say it once more.
\end{tcolorbox}
\section[\texorpdfstring{\sk{उज्ज्वलः} (ujjvalaḥ)}{उज्ज्वलः (ujjvalaḥ)}]{\sk{उज्ज्वलः} (ujjvalaḥ) — bright, shining}

\label{lesson:SA-C188-ujjvalah}



\begin{quote}
Before the new one: say the Sanskrit for busy, then the Sanskrit for spotless.
\end{quote}

\subsection*{You'll want to know: \sk{उज्ज्वलः}}
\textbf{\sk{उज्ज्वलः}} — \emph{ujjvalaḥ} — \textquotedblleft{}bright, shining\textquotedblright{}.

\begin{grammarlens}[title={What you've built}]
One more word for this chapter.
\end{grammarlens}

\subsection*{Your turn}
\begin{itemize}
\raggedright
  \item \emph{Say it:} \emph{ujjvalaḥ}
  \item \emph{Say it:} \emph{ujjvalaḥ}, once more
  \item \emph{Say it:} say \emph{ujjvalaḥ}
  \item \emph{Recall:} say the Sanskrit for busy, then the Sanskrit for spotless, then say \emph{ujjvalaḥ} again
\end{itemize}

\begin{tcolorbox}[breakable,colback=teal!4,colframe=teal!35!black,title={Before you move on}]
What does \sk{उज्ज्वलः} mean? (\textquotedblleft{}Bright, shining\textquotedblright{}.) Say it once more.
\end{tcolorbox}
